\documentclass[letterpaper,12pt]{article}

\usepackage{fancyhdr}
\usepackage[most]{tcolorbox}
\usepackage{longtable}
\usepackage{makecell}
\usepackage{tikz}
\usepackage{geometry}
\usepackage{array}
\usepackage{multicol}
\usepackage[table]{xcolor}
\usepackage{polyglossia}
\setmainlanguage{english}
\setotherlanguage{arabic}
\setotherlanguage{hebrew}
\usepackage{fontspec}
\setmainfont{Charis SIL}
\newfontfamily\arabicfont[Script=Arabic,Scale=1.2]{Amiri}
\newfontfamily\hebrewfont{Ezra SIL}[Script=Hebrew]
\geometry{margin=1.5cm}
\usepackage{booktabs}
\usepackage{graphicx}
\usepackage{multirow}

\newfontfamily\boxdrawfont{DejaVu Sans Mono}[Scale=0.9]

\setlength{\headheight}{14.49998pt}
\addtolength{\topmargin}{-2.49998pt}
\pagestyle{fancy}
\fancyhf{}
\fancyfoot[C]{\thepage}
\fancyhead[R]{\textit{What Use Is Speaking Words from Quraysh Times?}}
\renewcommand{\headrulewidth}{0pt}

\usetikzlibrary{shapes,arrows,decorations.pathmorphing}

% Custom colors
\definecolor{headercolor}{RGB}{70,130,180}
\definecolor{boxcolor}{RGB}{240,248,255}
\definecolor{accentcolor}{RGB}{220,20,60}
\definecolor{tableheader}{RGB}{220,220,220}
\definecolor{dialectcolor}{RGB}{34,139,34}

\begin{document}

\title{\textbf{\Large Arabic Phrase Analyser}\\
\large What Use Is Speaking Words from Quraysh Times?\\
\normalsize \textarabic{ما جدوى النطق بكلمات من زمن قريش محكية}}
\author{Igor Deruga}
\date{}
\maketitle

% ======================== Phrase Display ========================
\begin{tcolorbox}[colback=boxcolor,colframe=headercolor,title=\textbf{Arabic Phrase with Full Diacritics},breakable]
\centering
\Large\textarabic{مَا جَدْوَى النُّطْقِ بِكَلِمَاتٍ مِنْ زَمَنِ قُرَيْشٍ مُحَكِّيَةٍ؟}
\end{tcolorbox}

% ======================== Translations ========================
\section{English Translation}
\begin{tcolorbox}[colback=white,colframe=accentcolor,breakable]
\textbf{Literal:} What usefulness the-speaking with-words from time Quraysh spoken?\\
\textit{[Arabic order retained for direct mapping]}\\[0.5em]
\textbf{Adapted:} What is the use of speaking words from the time of Quraysh that are spoken?
\end{tcolorbox}

% ======================== Detailed Word Analysis ========================
\section{Detailed Word Analysis}

\subsection{\textarabic{مَا} — /maː/}
\begin{tabular}{p{3cm}p{10cm}}
\toprule
\textbf{Translation} & what \\
\textbf{Root} & \textarabic{م-ا} (interrogative particle, not from triliteral root) \\
\textbf{Pattern} & Independent interrogative particle \\
\textbf{Grammar} & Interrogative particle introducing questions about identity, nature, or quality \\
\midrule
\textbf{Examples} & \makecell[l]{\parbox{9.5cm}{
1. \textarabic{مَا اسْمُكَ؟} — What is your name? /maː sːmuka/\\
2. \textarabic{مَا هَذَا؟} — What is this? /maː haːðaː/\\
3. \textarabic{مَا الَّذِي حَدَثَ؟} — What happened? /maː llaːðiː ħadaθa/
}} \\
\midrule
\textbf{Synonyms} & \textarabic{مَاذَا} (what - with \textarabic{ذَا}), \textarabic{أَيّ} (which/what kind) \\
\textbf{Etymology} & From Proto-Semitic *mā (interrogative), related to Hebrew \texthebrew{מָה} (mah) "what" \\
\bottomrule
\end{tabular}

\subsubsection*{Usage Notes}
\begin{itemize}
  \item \textarabic{مَا} can function as interrogative "what" or as negation "not" (context-dependent)
  \item When interrogative, it can be followed by nouns or verbal sentences
  \item Can combine with other particles: \textarabic{مَاذَا} (maːðaː), \textarabic{لِمَا} (limaː "why")
  \item As negation particle (different usage): \textarabic{مَا قَالَ} "he did not say"
\end{itemize}

\subsection{\textarabic{جَدْوَى} — /dʒadwaː/}
\begin{tabular}{p{3cm}p{10cm}}
\toprule
\textbf{Translation} & usefulness, utility, benefit \\
\textbf{Root} & \textarabic{ج-د-و} (dʒ-d-w) \\
\textbf{Pattern} & \textarabic{فَعْلَى} (faʕlaː) — feminine noun pattern \\
\textbf{Grammar} & Feminine noun, nominative case (subject of nominal sentence), indefinite \\
\midrule
\textbf{Examples} & \makecell[l]{\parbox{9.5cm}{
1. \textarabic{لَا جَدْوَى مِنَ الْبُكَاءِ} — There is no use in crying /laː dʒadwaː mina lbukaːʔi/\\
2. \textarabic{مَا جَدْوَى هَذَا الْعَمَلِ؟} — What is the benefit of this work? /maː dʒadwaː haːða lʕamali/\\
3. \textarabic{لَيْسَ لَهُ جَدْوَى} — It has no usefulness /lajsa lahu dʒadwaː/
}} \\
\midrule
\textbf{Synonyms} & \textarabic{فَائِدَة} (benefit), \textarabic{نَفْع} (utility), \textarabic{جَدِير} (worthwhile) \\
\textbf{Etymology} & From root \textarabic{ج-د-و} meaning "to be beneficial, to avail" \\
\bottomrule
\end{tabular}

\subsubsection*{Declension}
\begin{tabular}{|c|c|c|}
\hline
\textbf{Case} & \textbf{Indefinite} & \textbf{Definite} \\
\hline
Nominative & \textarabic{جَدْوَى} /dʒadwaː/ & \textarabic{الْجَدْوَى} /aldʒadwaː/ \\
\hline
Accusative & \textarabic{جَدْوَى} /dʒadwaː/ & \textarabic{الْجَدْوَى} /aldʒadwaː/ \\
\hline
Genitive & \textarabic{جَدْوَى} /dʒadwaː/ & \textarabic{الْجَدْوَى} /aldʒadwaː/ \\
\hline
\end{tabular}
\vspace{0.5em}

\textit{Note: As a feminine singular noun ending in alif maqṣūrah (\textarabic{ى}), it remains unchanged across all three cases.}

\subsection{\textarabic{النُّطْقِ} — /nnuːtˤqi/}
\begin{tabular}{p{3cm}p{10cm}}
\toprule
\textbf{Translation} & the speaking, the pronunciation, the utterance \\
\textbf{Root} & \textarabic{ن-ط-ق} (n-tˤ-q) \\
\textbf{Pattern} & \textarabic{فُعْل} (fuʕl) — verbal noun (maṣdar) pattern \\
\textbf{Grammar} & Masculine noun, genitive case (governed by \textarabic{جَدْوَى} in iḍāfah construction), definite \\
\midrule
\textbf{Examples} & \makecell[l]{\parbox{9.5cm}{
1. \textarabic{النُّطْقُ الصَّحِيحُ مُهِمٌّ} — Correct pronunciation is important /annuːtˤqu sˤsˤaħiːħu muhimmun/\\
2. \textarabic{يَصْعُبُ النُّطْقُ بِهَذِهِ الْكَلِمَةِ} — This word is difficult to pronounce /jasʕubu nnuːtˤqu bihaːðihi lkalimati/\\
3. \textarabic{مَلَكَةُ النُّطْقِ} — The faculty of speech /malaakatu nnuːtˤqi/
}} \\
\midrule
\textbf{Synonyms} & \textarabic{كَلَام} (speech), \textarabic{تَلَفُّظ} (articulation), \textarabic{لَفْظ} (utterance) \\
\textbf{Etymology} & From root meaning "to speak, to utter, to articulate" \\
\bottomrule
\end{tabular}

\subsubsection*{Related Verbal Forms}
\par{\large \textbf{Verbal Noun (Maṣdar)}: \textarabic{نُطْق} /nutˤq/}

\begin{longtable}{|>{\raggedright}p{3.5cm}|p{5cm}|p{5cm}|}
\hline
\textbf{Person} & \textbf{Perfect (Past)} & \textbf{Imperfect (Present)} \\
\hline
\textbf{3rd person masc. sing.} & \textarabic{نَطَقَ} /natˤaqa/ & \textarabic{يَنْطِقُ} /jantˤiqu/ \\
\hline
\textbf{3rd person fem. sing.} & \textarabic{نَطَقَتْ} /natˤaqat/ & \textarabic{تَنْطِقُ} /tantˤiqu/ \\
\hline
\textbf{3rd person masc. plural} & \textarabic{نَطَقُوا} /natˤaquː/ & \textarabic{يَنْطِقُونَ} /jantˤiquːna/ \\
\hline
\textbf{2nd person masc. sing.} & \textarabic{نَطَقْتَ} /natˤaqta/ & \textarabic{تَنْطِقُ} /tantˤiqu/ \\
\hline
\textbf{1st person singular} & \textarabic{نَطَقْتُ} /natˤaqtu/ & \textarabic{أَنْطِقُ} /ʔantˤiqu/ \\
\hline
\end{longtable}

\subsection{\textarabic{بِكَلِمَاتٍ} — /bikalimaːtin/}
\begin{tabular}{p{3cm}p{10cm}}
\toprule
\textbf{Translation} & with words \\
\textbf{Root} & \textarabic{ك-ل-م} (k-l-m) \\
\textbf{Pattern} & \textarabic{فَعِلَة} (faʕilah) — plural: \textarabic{فَعِلَات} (faʕilaːt) \\
\textbf{Grammar} & Preposition \textarabic{بِ} + feminine noun plural, genitive case (after preposition), indefinite with tanwīn \\
\midrule
\textbf{Examples} & \makecell[l]{\parbox{9.5cm}{
1. \textarabic{كَلِمَةٌ وَاحِدَةٌ} — One word /kalimatun waːħidatun/\\
2. \textarabic{تَكَلَّمَ بِكَلِمَاتٍ جَمِيلَةٍ} — He spoke with beautiful words /takallama bikalimaːtin dʒamiːlatin/\\
3. \textarabic{الْكَلِمَاتُ الثَّلَاثُ} — The three words /alkalimaːtu θθalaːθu/
}} \\
\midrule
\textbf{Synonyms} & \textarabic{أَلْفَاظ} (words/expressions), \textarabic{عِبَارَات} (phrases), \textarabic{مُفْرَدَات} (vocabulary) \\
\textbf{Etymology} & From root meaning "to speak, to wound"; related to Hebrew \texthebrew{כָּלַם} "to humiliate, wound" \\
\bottomrule
\end{tabular}

\subsubsection*{Declension of \textarabic{كَلِمَة}}
\begin{tabular}{|c|c|c|c|}
\hline
\textbf{Number} & \textbf{Case} & \textbf{Indefinite} & \textbf{Definite} \\
\hline
\multirow{3}{*}{Singular} & Nominative & \textarabic{كَلِمَةٌ} /kalimatun/ & \textarabic{الْكَلِمَةُ} /alkalimatu/ \\
& Accusative & \textarabic{كَلِمَةً} /kalimatan/ & \textarabic{الْكَلِمَةَ} /alkalimata/ \\
& Genitive & \textarabic{كَلِمَةٍ} /kalimatin/ & \textarabic{الْكَلِمَةِ} /alkalimati/ \\
\hline
\multirow{3}{*}{Plural} & Nominative & \textarabic{كَلِمَاتٌ} /kalimaːtun/ & \textarabic{الْكَلِمَاتُ} /alkalimaːtu/ \\
& Accusative & \textarabic{كَلِمَاتٍ} /kalimaːtin/ & \textarabic{الْكَلِمَاتِ} /alkalimaːti/ \\
& Genitive & \textarabic{كَلِمَاتٍ} /kalimaːtin/ & \textarabic{الْكَلِمَاتِ} /alkalimaːti/ \\
\hline
\end{tabular}

\subsection{\textarabic{مِنْ} — /min/}
\begin{tabular}{p{3cm}p{10cm}}
\toprule
\textbf{Translation} & from \\
\textbf{Root} & Preposition (not from triliteral root) \\
\textbf{Pattern} & Independent preposition \\
\textbf{Grammar} & Preposition indicating source, origin, or partitive meaning \\
\midrule
\textbf{Examples} & \makecell[l]{\parbox{9.5cm}{
1. \textarabic{جِئْتُ مِنَ الْبَيْتِ} — I came from the house /dʒiʔtu mina lbajti/\\
2. \textarabic{هَذَا مِنْ فَضْلِ اللهِ} — This is from God's bounty /haːðaː min fadˤli llaːhi/\\
3. \textarabic{وَاحِدٌ مِنَّا} — One of us /waːħidun minnaː/
}} \\
\midrule
\textbf{Synonyms} & (context-dependent): \textarabic{عَنْ} (from/about), \textarabic{عِنْدَ} (from/at) \\
\textbf{Etymology} & From Proto-Semitic *min; Hebrew \texthebrew{מִן} (min) "from" \\
\bottomrule
\end{tabular}

\subsection{\textarabic{زَمَنِ} — /zamani/}
\begin{tabular}{p{3cm}p{10cm}}
\toprule
\textbf{Translation} & time (of) \\
\textbf{Root} & \textarabic{ز-م-ن} (z-m-n) \\
\textbf{Pattern} & \textarabic{فَعَل} (faʕal) \\
\textbf{Grammar} & Masculine noun, genitive case (after preposition \textarabic{مِنْ}), first term of iḍāfah (construct state) \\
\midrule
\textbf{Examples} & \makecell[l]{\parbox{9.5cm}{
1. \textarabic{فِي ذَلِكَ الزَّمَانِ} — In that time /fiː ðaːlika zzamaːni/\\
2. \textarabic{مَرَّ زَمَنٌ طَوِيلٌ} — A long time passed /marra zamanun tˤawiːlun/\\
3. \textarabic{زَمَنُ الْعَرَبِ} — The time of the Arabs /zamanu lʕarabi/
}} \\
\midrule
\textbf{Synonyms} & \textarabic{وَقْت} (time), \textarabic{عَصْر} (era/age), \textarabic{دَهْر} (age/epoch) \\
\textbf{Etymology} & From root meaning "time, era, period" \\
\bottomrule
\end{tabular}

\subsubsection*{Declension}
\begin{tabular}{|c|c|c|}
\hline
\textbf{Case} & \textbf{Indefinite} & \textbf{Definite} \\
\hline
Nominative & \textarabic{زَمَنٌ} /zamanun/ & \textarabic{الزَّمَنُ} /azzamanu/ \\
\hline
Accusative & \textarabic{زَمَنًا} /zamanan/ & \textarabic{الزَّمَنَ} /azzamana/ \\
\hline
Genitive & \textarabic{زَمَنٍ} /zamanin/ & \textarabic{الزَّمَنِ} /azzamani/ \\
\hline
\end{tabular}

\subsection{\textarabic{قُرَيْشٍ} — /qurajiʃin/}
\begin{tabular}{p{3cm}p{10cm}}
\toprule
\textbf{Translation} & Quraysh (the dominant Arabian tribe at the time of Prophet Muhammad) \\
\textbf{Root} & \textarabic{ق-ر-ش} (q-r-ʃ) \\
\textbf{Pattern} & \textarabic{فُعَيْل} (fuʕajl) — diminutive pattern \\
\textbf{Grammar} & Proper noun (tribe name), feminine, genitive case (second term of iḍāfah), diptote (no tanwīn) \\
\midrule
\textbf{Examples} & \makecell[l]{\parbox{9.5cm}{
1. \textarabic{قُرَيْشٌ قَبِيلَةٌ عَرَبِيَّةٌ} — Quraysh is an Arab tribe /qurajiʃun qabiːlatun ʕarabijjatun/\\
2. \textarabic{لُغَةُ قُرَيْشٍ} — The language of Quraysh /luɣatu qurajiʃin/\\
3. \textarabic{كَانَتْ قُرَيْشٌ تَسْكُنُ مَكَّةَ} — Quraysh used to inhabit Mecca /kaːnat qurajiʃun taskunu makkah/
}} \\
\midrule
\textbf{Cultural Note} & The Quraysh tribe controlled Mecca and the Kaaba. Prophet Muhammad was from Quraysh. The Quranic Arabic is based on the Qurayshi dialect, considered the purest form of Arabic. \\
\textbf{Etymology} & Possibly from root meaning "to gather, collect" or diminutive form; various etymological theories exist \\
\bottomrule
\end{tabular}

\subsection{\textarabic{مُحَكِّيَةٍ} — /muħakkijatin/}
\begin{tabular}{p{3cm}p{10cm}}
\toprule
\textbf{Translation} & spoken, narrated, recounted \\
\textbf{Root} & \textarabic{ح-ك-ي} (ħ-k-j) \\
\textbf{Pattern} & \textarabic{مُفَعِّلَة} (mufaʕʕilah) — Form II active participle, feminine \\
\textbf{Grammar} & Active participle, feminine plural genitive (agreeing with \textarabic{كَلِمَاتٍ}), indefinite with tanwīn \\
\midrule
\textbf{Examples} & \makecell[l]{\parbox{9.5cm}{
1. \textarabic{حَكَى لِي قِصَّةً} — He narrated a story to me /ħakaː liː qisˤsˤatan/\\
2. \textarabic{يُحَكِّي التَّارِيخَ} — He recounts history /juħakkiː ttaːriːxa/\\
3. \textarabic{الْقِصَّةُ الْمُحَكِّيَّةُ} — The narrated story /alqisˤsˤatu lmuħakkijjatu/
}} \\
\midrule
\textbf{Synonyms} & \textarabic{مَرْوِيَّة} (narrated), \textarabic{مَنْقُولَة} (transmitted), \textarabic{مُتَحَدَّث بِهَا} (spoken) \\
\textbf{Etymology} & From root meaning "to tell, narrate, imitate, speak" \\
\bottomrule
\end{tabular}

\subsubsection*{Form II Conjugation of \textarabic{حَكَّى}}
\par{\large \textbf{Verbal Noun}: \textarabic{تَحْكِيَة} /taħkijjah/}

\begin{longtable}{|>{\raggedright}p{3.5cm}|p{5cm}|p{5cm}|}
\hline
\textbf{Person} & \textbf{Perfect (Past)} & \textbf{Imperfect (Present)} \\
\hline
\textbf{3rd person masc. sing.} & \textarabic{حَكَّى} /ħakkaː/ & \textarabic{يُحَكِّي} /juħakkiː/ \\
\hline
\textbf{3rd person fem. sing.} & \textarabic{حَكَّتْ} /ħakkat/ & \textarabic{تُحَكِّي} /tuħakkiː/ \\
\hline
\textbf{3rd person masc. plural} & \textarabic{حَكَّوْا} /ħakkaw/ & \textarabic{يُحَكُّونَ} /juħakkuːna/ \\
\hline
\textbf{2nd person masc. sing.} & \textarabic{حَكَّيْتَ} /ħakkajta/ & \textarabic{تُحَكِّي} /tuħakkiː/ \\
\hline
\textbf{1st person singular} & \textarabic{حَكَّيْتُ} /ħakkajtu/ & \textarabic{أُحَكِّي} /ʔuħakkiː/ \\
\hline
\end{longtable}

\subsubsection*{Active Participle Forms}
\begin{tabular}{|c|c|c|}
\hline
\textbf{Gender/Number} & \textbf{Indefinite} & \textbf{Definite} \\
\hline
Masculine singular & \textarabic{مُحَكٍّ} /muħakkin/ & \textarabic{الْمُحَكِّي} /almuħakkiː/ \\
\hline
Feminine singular & \textarabic{مُحَكِّيَةٌ} /muħakkijatun/ & \textarabic{الْمُحَكِّيَةُ} /almuħakkijatu/ \\
\hline
Masculine plural & \textarabic{مُحَكُّونَ} /muħakkuːna/ & \textarabic{الْمُحَكُّونَ} /almuħakkuːna/ \\
\hline
Feminine plural & \textarabic{مُحَكِّيَاتٌ} /muħakkijaːtun/ & \textarabic{الْمُحَكِّيَاتُ} /almuħakkijaːtu/ \\
\hline
\end{tabular}

% ======================== Phrase Analysis ========================
\section{Phrase Analysis}

\begin{tcolorbox}[colback=boxcolor,colframe=headercolor,breakable]
\textbf{Grammatical Structure:}\\
Interrogative particle + noun (subject) + definite verbal noun (genitive in iḍāfah) + preposition + noun plural (genitive) + preposition + noun (genitive in iḍāfah) + proper noun (genitive) + active participle (adjective, genitive)

\vspace{0.5em}
\textbf{Sentence Type:} Nominal sentence (\textarabic{جُمْلَة اسْمِيَّة}) introduced by interrogative \textarabic{مَا}

\vspace{0.5em}
\textbf{Key Grammar Points:}
\begin{itemize}
\item \textbf{Interrogative \textarabic{مَا}:} Opens rhetorical question about utility/benefit
\item \textbf{Iḍāfah construction 1:} \textarabic{جَدْوَى النُّطْقِ} (the usefulness of speaking) — possessive/genitive relationship
\item \textbf{Prepositional phrase 1:} \textarabic{بِكَلِمَاتٍ} (with words) — describes the manner of speaking
\item \textbf{Iḍāfah construction 2:} \textarabic{زَمَنِ قُرَيْشٍ} (the time of Quraysh) — temporal-possessive relationship
\item \textbf{Adjectival phrase:} \textarabic{مُحَكِّيَةٍ} (spoken) — Form II active participle modifying \textarabic{كَلِمَاتٍ}
\item \textbf{Agreement:} The participle \textarabic{مُحَكِّيَةٍ} agrees with \textarabic{كَلِمَاتٍ} in gender (feminine), number (plural), case (genitive), and definiteness (indefinite)
\item \textbf{Rhetorical nature:} Question implies skepticism about the value of using archaic Qurayshi Arabic in modern contexts
\item \textbf{Diptote proper noun:} \textarabic{قُرَيْشٍ} lacks tanwīn in genitive case (characteristic of many proper nouns and certain patterns)
\end{itemize}
\end{tcolorbox}

\section{Syntactic Analysis
}

\begin{tcolorbox}[colback=white,colframe=accentcolor,breakable]
\textbf{Phrase Structure Tree:}

\begin{center}
\begin{tikzpicture}[
  box/.style={rectangle, draw, minimum width=2.5cm, minimum height=0.8cm, align=center},
  arrow/.style={->, >=stealth, thick}
]

% Main question word
\node[box] (ma) at (0,0) {\textarabic{مَا}\\What?};

% Main predicate
\node[box, minimum width=8cm] (pred) at (6,0) {What is being questioned};

% Subject idafa
\node[box] (jadwa) at (3,-2) {\textarabic{جَدْوَى}\\usefulness};
\node[box] (nutq) at (6,-2) {\textarabic{النُّطْقِ}\\the speaking};

% Prepositional phrases
\node[box] (bikalimat) at (2,-4) {\textarabic{بِكَلِمَاتٍ}\\with words};
\node[box] (min) at (6,-4) {\textarabic{مِنْ}\\from};
\node[box] (zaman) at (9,-4) {\textarabic{زَمَنِ}\\time};

% Bottom level
\node[box] (muhakkiya) at (2,-6) {\textarabic{مُحَكِّيَةٍ}\\spoken};
\node[box] (quraysh) at (9,-6) {\textarabic{قُرَيْشٍ}\\Quraysh};

% Arrows
\draw[arrow] (ma) -- (pred);
\draw[arrow] (pred) -- (jadwa);
\draw[arrow] (jadwa) -- (nutq);
\draw[arrow] (nutq) -- (bikalimat);
\draw[arrow] (nutq) -- (min);
\draw[arrow] (bikalimat) -- (muhakkiya);
\draw[arrow] (min) -- (zaman);
\draw[arrow] (zaman) -- (quraysh);

\end{tikzpicture}
\end{center}

\textbf{Syntactic Relationships:}
\begin{enumerate}
\item \textbf{Main clause:} \textarabic{مَا جَدْوَى...} — Interrogative nominal sentence
\item \textbf{First iḍāfah:} \textarabic{جَدْوَى النُّطْقِ} — "usefulness of speaking" (muḍāf + muḍāf ilayh)
\item \textbf{Adjectival modification:} \textarabic{النُّطْقِ} is further specified by prepositional phrase and relative construction
\item \textbf{Prepositional phrase:} \textarabic{بِكَلِمَاتٍ} — instrument/means of speaking
\item \textbf{Second iḍāfah:} \textarabic{مِنْ زَمَنِ قُرَيْشٍ} — temporal origin phrase
\item \textbf{Participial adjective:} \textarabic{مُحَكِّيَةٍ} — modifies \textarabic{كَلِمَاتٍ}, shows agreement in all grammatical categories
\end{enumerate}
\end{tcolorbox}

% ======================== Similar Phrases for Practice ========================
\section{Similar Phrases for Practice}

\begin{enumerate}
\item \textarabic{مَا جَدْوَى الْكِتَابَةِ بِلُغَةٍ مِنْ عَصْرِ الْجَاهِلِيَّةِ مَنْسِيَّةٍ؟}\\
What is the use of writing in a language from the pre-Islamic era that is forgotten? /maː dʒadwaː lkitaːbati biluɣatin min ʕasˤri ldʒaːhilijjati mansijjatin/

\item \textarabic{مَا فَائِدَةُ الْحَدِيثِ بِعِبَارَاتٍ مِنْ زَمَنِ الْأَجْدَادِ مَهْجُورَةٍ؟}\\
What is the benefit of speaking with expressions from the time of ancestors that are abandoned? /maː faːʔidatu lħadiːθi biʕibaːraːtin min zamani l-ʔadʒdaːdi mahdʒuːratin/

\item \textarabic{مَا نَفْعُ الْقِرَاءَةِ بِكَلِمَاتٍ مِنْ تَارِيخِ الْعَرَبِ مَتْرُوكَةٍ؟}\\
What is the utility of reading with words from Arab history that are discarded? /maː nafʕu lqiraːʔati bikalimaːtin min taːriːxi lʕarabi matruːkatin/

\item \textarabic{مَا الْفَائِدَةُ مِنَ التَّكَلُّمِ بِأَلْفَاظٍ مِنْ لُغَةِ قُرَيْشٍ قَدِيمَةٍ؟}\\
What is the benefit of speaking with words from the language of Quraysh that are ancient? /maː lfaːʔidatu mina ttakallumi bi-ʔalfaːðˤin min luɣati qurajiʃin qadiːmatin/

\item \textarabic{مَا جَدْوَى الْحِوَارِ بِمُفْرَدَاتٍ مِنْ عَهْدِ النَّبِيِّ مَنْقُولَةٍ؟}\\
What is the use of dialogue with vocabulary from the Prophet's era that is transmitted? /maː dʒadwaː lħiwaːri bimufradaːtin min ʕahdi nnabijji manquːlatin/

\item \textarabic{هَلْ هُنَاكَ جَدْوَى مِنَ اسْتِعْمَالِ كَلِمَاتٍ مِنْ زَمَنٍ مَاضٍ؟}\\
Is there any use in using words from a past time? /hal hunaːka dʒadwaː mina stiʕmaːli kalimaːtin min zamanin maːðˤin/
\end{enumerate}

% ======================== Levantine Dialect Version ========================
\section{Levantine (Shami) Arabic Dialect}

\begin{tcolorbox}[colback=white,colframe=dialectcolor,title=\textbf{Levantine Version},breakable]
\textarabic{شُو فَايْدِة الْحَكِي بِكِلْمَات مِنْ زَمَان قُرَيْش مَحْكِيِّي؟}\\[0.5em]
\textbf{Phonetic:} /ʃuː faːjde lħakiː biklmaːt min zamaːn qurajiʃ maħkijjiː/\\[0.5em]
\textbf{Translation:} What's the benefit of speaking with words from Quraysh times that are spoken?
\end{tcolorbox}

\textbf{Key Dialectal Changes:}
\begin{itemize}
\item \textarabic{مَا} → \textarabic{شُو} /ʃuː/ — Standard Levantine interrogative "what"
\item \textarabic{جَدْوَى} → \textarabic{فَايْدِة} /faːjde/ — More common dialectal word for "benefit/use"
\item \textarabic{النُّطْق} → \textarabic{الْحَكِي} /lħakiː/ — Dialectal verbal noun "speaking/talking"
\item \textarabic{بِكَلِمَاتٍ} → \textarabic{بِكِلْمَات} /biklmaːt/ — Loss of case endings and short vowels
\item \textarabic{مِنْ زَمَنِ} → \textarabic{مِنْ زَمَان} /min zamaːn/ — Dialectal form (elongated final vowel)
\item \textarabic{مُحَكِّيَةٍ} → \textarabic{مَحْكِيِّي} /maħkijjiː/ — Simplified participle form without case endings
\item Overall simplification: Loss of iʕrāb (case endings), tanwīn, and some short vowels
\item Retention of meaning but with colloquial vocabulary and phonology
\end{itemize}

\subsection{Alternative Dialectal Versions}

\textbf{Egyptian Arabic:}\\
\textarabic{إِيهِ الْفَايْدَة مِنِ الْكَلَام بِكَلِمَات مِنْ زَمَان قُرَيْش مِتْقَالَة؟}\\
/ʔeː elfaːjda meni lkalaːm bekalimaːt min zamaːn qurajiʃ metʔaːla/

\textbf{Gulf Arabic:}\\
\textarabic{شُو الْفَايْدَة مِنْ الْحَچِي بِكَلِمَات مِنْ أَيَّام قُرَيْش مَحْچِيَّة؟}\\
/ʃuː elfaːjda min elħatʃiː bekalimaːt min ajjaːm qurajiʃ maħtʃijja/

% ======================== Additional Learning Notes ========================
\section{Additional Learning Notes}

\begin{tcolorbox}[colback=boxcolor,colframe=accentcolor,title=\textbf{Cultural and Linguistic Context},breakable]

\textbf{Historical Significance of Qurayshi Arabic:}
\begin{itemize}
\item The Quraysh tribe was the custodian of the Kaaba in pre-Islamic Mecca and the most prestigious Arab tribe
\item Their dialect became the foundation of Classical Arabic and the language of the Quran
\item Qurayshi Arabic was considered the purest and most eloquent form of Arabic in the 7th century
\item The phrase questions the relevance of using archaic Classical Arabic forms in modern contexts
\end{itemize}

\textbf{Linguistic Debate:}
This phrase touches on ongoing debates in the Arab world:
\begin{itemize}
\item \textbf{Fuṣḥā vs. ʕĀmmiyyah:} Standard/Classical Arabic vs. colloquial dialects
\item \textbf{Language modernization:} Should Arabic evolve or preserve classical forms?
\item \textbf{Diglossia:} The coexistence of formal and informal language registers
\item \textbf{Educational policy:} What form of Arabic should be taught and used?
\end{itemize}

\textbf{Rhetorical Function:}
The question is typically rhetorical, implying one of two positions:
\begin{enumerate}
\item \textit{Skeptical view:} Ancient forms are impractical for modern communication
\item \textit{Defensive view:} Classical forms maintain linguistic heritage and precision
\end{enumerate}

\textbf{Grammatical Complexity:}
This phrase demonstrates several advanced Arabic features:
\begin{itemize}
\item Multiple iḍāfah constructions (possessive chains)
\item Active participle as adjective with full agreement
\item Proper noun declension (diptote pattern)
\item Layered prepositional and adjectival modification
\end{itemize}
\end{tcolorbox}

\subsection{Memory Tips}

\begin{enumerate}
\item \textbf{Question Structure:} Remember the pattern \textarabic{مَا جَدْوَى...؟} for asking "What's the use of...?" — commonly used in rhetorical questions

\item \textbf{Iḍāfah Chain:} Visualize: \textit{usefulness} ← \textit{of speaking} ← \textit{with words} ← \textit{from time} ← \textit{of Quraysh}\\
Each element specifies the previous one more precisely

\item \textbf{Participle Agreement:} \textarabic{مُحَكِّيَةٍ} matches \textarabic{كَلِمَاتٍ} in four ways:
\begin{itemize}
\item Gender: feminine (\textarabic{ـَة})
\item Number: plural (sound feminine plural pattern)
\item Case: genitive (\textarabic{ـِ})
\item Definiteness: indefinite (tanwīn \textarabic{ـٍ})
\end{itemize}

\item \textbf{Root Recognition:} 
\begin{itemize}
\item \textarabic{ن-ط-ق} → speaking/pronunciation
\item \textarabic{ح-ك-ي} → narrating/telling/speaking
\item \textarabic{ك-ل-م} → words/speech
\end{itemize}

\item \textbf{Diptote Alert:} Proper nouns like \textarabic{قُرَيْش} often lack tanwīn — watch for this pattern with tribal names and place names

\item \textbf{Mnemonic for \textarabic{جَدْوَى}:} "Jadwa" sounds like "Java" (coffee) → "What's the use of coffee?" → asking about usefulness/benefit
\end{enumerate}

\subsection{Related Vocabulary Families}

\begin{multicols}{2}

\textbf{Usefulness/Benefit:}\\
\textarabic{جَدْوَى} — usefulness /dʒadwaː/\\
\textarabic{فَائِدَة} — benefit /faːʔida/\\
\textarabic{نَفْع} — utility /nafʕ/\\
\textarabic{مَنْفَعَة} — advantage /manfaʕa/\\
\textarabic{مَصْلَحَة} — interest /masˤlaħa/\\
\textarabic{غَايَة} — purpose /ɣaːja/\\

\textbf{Speech/Speaking:}\\
\textarabic{نُطْق} — pronunciation /nutˤq/\\
\textarabic{كَلَام} — speech /kalaːm/\\
\textarabic{حَدِيث} — talk /ħadiːθ/\\
\textarabic{قَوْل} — saying /qawl/\\
\textarabic{لَفْظ} — utterance /lafðˤ/\\
\textarabic{تَعْبِير} — expression /taʕbiːr/\\

\textbf{Time/Era:}\\
\textarabic{زَمَن} — time /zaman/\\
\textarabic{وَقْت} — time /waqt/\\
\textarabic{عَصْر} — era /ʕasˤr/\\
\textarabic{عَهْد} — epoch /ʕahd/\\
\textarabic{دَهْر} — age /dahr/\\
\textarabic{حِقْبَة} — period /ħiqba/\\

\textbf{Words/Language:}\\
\textarabic{كَلِمَة} — word /kalima/\\
\textarabic{لَفْظَة} — word /lafðˤa/\\
\textarabic{مُفْرَدَة} — vocabulary item /mufrada/\\
\textarabic{عِبَارَة} — phrase /ʕibaːra/\\
\textarabic{لُغَة} — language /luɣa/\\
\textarabic{لِسَان} — tongue/language /lisaːn/\\

\textbf{Narration/Speaking:}\\
\textarabic{حِكَايَة} — narration /ħikaːja/\\
\textarabic{رِوَايَة} — narrative /riwaːja/\\
\textarabic{قِصَّة} — story /qisˤsˤa/\\
\textarabic{سَرْد} — narration /sard/\\
\textarabic{إِخْبَار} — telling /ʔixbaːr/\\
\textarabic{نَقْل} — transmission /naql/\\

\textbf{Adjectives (archaic/old):}\\
\textarabic{قَدِيم} — ancient /qadiːm/\\
\textarabic{عَتِيق} — old /ʕatiːq/\\
\textarabic{مَاضٍ} — past /maːðˤin/\\
\textarabic{بَالٍ} — worn out /baːlin/\\
\textarabic{مَنْسِيّ} — forgotten /mansijj/\\
\textarabic{مَهْجُور} — abandoned /mahdʒuːr/\\
\end{multicols}

\section{Grammatical Deep Dive}

\subsection{The Iḍāfah Construction (\textarabic{الْإِضَافَة})}

The phrase contains two iḍāfah constructions, demonstrating this fundamental Arabic grammatical structure:

\begin{tcolorbox}[colback=boxcolor,colframe=headercolor,breakable]
\textbf{Iḍāfah 1:} \textarabic{جَدْوَى النُّطْقِ} /dʒadwaː nnuːtˤqi/

\begin{tabular}{|l|l|l|}
\hline
\textbf{Component} & \textbf{Grammatical Role} & \textbf{Case} \\
\hline
\textarabic{جَدْوَى} (muḍāf) & Possessed noun (first term) & Nominative (subject) \\
\textarabic{النُّطْقِ} (muḍāf ilayh) & Possessing noun (second term) & Genitive (required) \\
\hline
\end{tabular}

\vspace{0.5em}
\textbf{Rules applied:}
\begin{itemize}
\item First term (\textarabic{جَدْوَى}) cannot take the definite article or tanwīn
\item First term takes its case from its syntactic function (here: nominative as subject)
\item Second term (\textarabic{النُّطْقِ}) must be genitive
\item Second term can be definite (making whole construction definite) or indefinite
\end{itemize}

\vspace{1em}
\textbf{Iḍāfah 2:} \textarabic{زَمَنِ قُرَيْشٍ} /zamani qurajiʃin/

\begin{tabular}{|l|l|l|}
\hline
\textbf{Component} & \textbf{Grammatical Role} & \textbf{Case} \\
\hline
\textarabic{زَمَنِ} (muḍāf) & Possessed noun & Genitive (after \textarabic{مِنْ}) \\
\textarabic{قُرَيْشٍ} (muḍāf ilayh) & Possessing noun & Genitive (required) \\
\hline
\end{tabular}

\vspace{0.5em}
\textbf{Special feature:}
\begin{itemize}
\item \textarabic{قُرَيْشٍ} is a diptote (\textarabic{مَمْنُوع مِنَ الصَّرْف}) — proper noun that doesn't take tanwīn
\item In genitive case, diptotes take kasrah without tanwīn: \textarabic{قُرَيْشٍ} not \textarabic{*قُرَيْشٍ}
\end{itemize}
\end{tcolorbox}

\subsection{Active Participle as Adjective}

\begin{tcolorbox}[colback=white,colframe=accentcolor,breakable]
\textbf{Participle:} \textarabic{مُحَكِّيَةٍ} /muħakkijatin/ — Form II active participle

\textbf{Base form (masc. sing.):} \textarabic{مُحَكٍّ} /muħakkin/

\textbf{Agreement Chain:}
\begin{center}
\begin{tabular}{|l|c|c|}
\hline
\textbf{Feature} & \textbf{\textarabic{كَلِمَاتٍ}} & \textbf{\textarabic{مُحَكِّيَةٍ}} \\
\hline
Gender & Feminine & Feminine (\textarabic{ـَة}) \\
Number & Plural & Plural (sound fem. pl.) \\
Case & Genitive & Genitive (\textarabic{ـِ}) \\
Definiteness & Indefinite & Indefinite (tanwīn \textarabic{ـٍ}) \\
\hline
\end{tabular}
\end{center}

\textbf{Why Form II?}\\
Form II (\textarabic{فَعَّلَ}) often indicates causative, intensive, or transitive meanings:
\begin{itemize}
\item Form I: \textarabic{حَكَى} — "to tell, narrate" (basic meaning)
\item Form II: \textarabic{حَكَّى} — "to narrate extensively, to imitate speech"
\item The active participle \textarabic{مُحَكِّيَة} → "that which is being narrated/spoken"
\end{itemize}
\end{tcolorbox}

\subsection{Rhetorical Question Formation}

\begin{tcolorbox}[colback=boxcolor,colframe=headercolor,breakable]
\textbf{Structure of \textarabic{مَا} questions:}

\begin{enumerate}
\item \textbf{Simple identification:}\\
\textarabic{مَا هَذَا؟} — What is this? (straightforward question)

\item \textbf{Quality/utility (rhetorical):}\\
\textarabic{مَا جَدْوَى...؟} — What's the use of...? (implies skepticism)\\
\textarabic{مَا فَائِدَةُ...؟} — What's the benefit of...? (questions value)

\item \textbf{Response patterns:}\\
Literal answer: \textarabic{جَدْوَاهُ كَذَا وَكَذَا} — Its use is such-and-such\\
Rhetorical response: \textarabic{لَا جَدْوَى مِنْهُ} — There's no use in it
\end{enumerate}

\textbf{Other rhetorical question patterns:}
\begin{itemize}
\item \textarabic{مَا الْفَائِدَةُ مِنْ...؟} — What's the benefit from...?
\item \textarabic{مَا النَّفْعُ فِي...؟} — What's the utility in...?
\item \textarabic{هَلْ هُنَاكَ جَدْوَى...؟} — Is there any use...?
\item \textarabic{أَيُّ فَائِدَةٍ...؟} — What benefit...?
\end{itemize}
\end{tcolorbox}

\section{Practice Exercises}

\subsection{Translation Practice}

Translate the following into Arabic using the same grammatical structures:

\begin{enumerate}
\item What is the use of reading books from ancient times that are abandoned?
\vspace{2em}

\item What is the benefit of writing letters from the past era that are forgotten?
\vspace{2em}

\item What is the utility of studying words from old Arabic that are unused?
\vspace{2em}
\end{enumerate}

\subsection{Grammatical Analysis}

Identify the following in the original phrase:

\begin{enumerate}
\item Circle all words in the genitive case
\item Underline the two iḍāfah constructions
\item Mark the active participle and show its agreement with its noun
\item Identify the interrogative particle
\end{enumerate}

\subsection{Transformation Exercises}

\begin{enumerate}
\item Change the phrase to past tense (if applicable): \textarabic{كَانَ مَا جَدْوَى...}
\vspace{1.5em}

\item Make it a statement instead of question: \textarabic{لَا جَدْوَى...}
\vspace{1.5em}

\item Change to Egyptian dialect
\vspace{1.5em}

\item Replace \textarabic{قُرَيْش} with \textarabic{الْجَاهِلِيَّة} and adjust grammar
\vspace{1.5em}
\end{enumerate}

\section{Answers to Practice Exercises}

\subsection{Translation Practice — Suggested Answers}

\begin{enumerate}
\item \textarabic{مَا جَدْوَى قِرَاءَةِ كُتُبٍ مِنْ أَزْمِنَةٍ قَدِيمَةٍ مَهْجُورَةٍ؟}\\
/maː dʒadwaː qiraːʔati kutubin min ʔazminatin qadiːmatin mahdʒuːratin/

\item \textarabic{مَا فَائِدَةُ كِتَابَةِ رَسَائِلَ مِنْ عَصْرٍ مَاضٍ مَنْسِيَّةٍ؟}\\
/maː faːʔidatu kitaːbati rasaːʔila min ʕasˤrin maːðˤin mansijjatin/

\item \textarabic{مَا نَفْعُ دِرَاسَةِ كَلِمَاتٍ مِنْ عَرَبِيَّةٍ عَتِيقَةٍ غَيْرِ مُسْتَعْمَلَةٍ؟}\\
/maː nafʕu diraːsati kalimaːtin min ʕarabijjatin ʕatiːqatin ɣajri mustaʕmalatin/
\end{enumerate}

\end{document}